\section*{Bab \ref{ch:b3:structural-analysis} --- Analisis Struktur}

\begin{solution}{exo:b3:structural-analysis:1}
\omterm{def:b3:structural-analysis:stereo}{Kutub} berarah $40 + 180 = \ang{220}$ dan menunjam $90 - 30 = \ang{60}$.
Bidang itu memiliki kemiringan $90 - 20 = \ang{70}$ ke arah $315 - 180 = \ang{135}$: $135/70$.
\end{solution}

\begin{solution}{exo:b3:structural-analysis:2}
$\theta = \ang{60}$. Sama sudut: $10\tan\ang{30} = \qty{5.77}{cm}$;
sama luas: $10\sqrt2\sin\ang{30} = \qty{7.07}{cm}$.
\end{solution}

\begin{solution}{exo:b3:structural-analysis:3}
$\vec n_1 = (\sin\ang{30}, 0, \cos\ang{30})$, $\vec n_2 = (-\sin\ang{30}, 0,
\cos\ang{30})$, $\vec n_1\cdot\vec n_2 = -0.25 + 0.75 = 0.5$: \omterm{def:b3:structural-analysis:stereo}{kutub-kutubnya}
terpisah sejauh \ang{60}, sudut antar-sayapnya \ang{120}, sebuah lipatan terbuka.
\end{solution}

\begin{solution}{exo:b3:structural-analysis:4}
$\vec d = (\cos\ang{60}, 0, -\sin\ang{60}) = (0.5, 0, -0.866)$ dan $\vec u =
0.707\,\vec s + 0.707\,\vec d = (0.354, 0.707, -0.612)$: tren
$\operatorname{atan2}(0.354, 0.707) = \ang{26.6}$, penunjaman $\arcsin0.612 =
\ang{37.8}$. Rake berada di antara \ang{-90} (normal) dan 0 (sinistral): sebuah
sesar normal miring dengan komponen sinistral.
\end{solution}

\begin{solution}{exo:b3:structural-analysis:5}
\begin{gather*}
  \vec n_1 = (0.5567, 0.3214, 0.7660),\qquad \vec n_2 = (-0.4698, -0.1710, 0.8660),\\
  \vec n_1\times\vec n_2 = (0.4093, -0.8420, 0.0558),
\end{gather*}
dengan panjang $0.9379$. Lawannya, $-\vec n_1\times\vec n_2$, menunjuk ke bawah:
tren \ang{334.1}, penunjaman $\arcsin(0.0558/0.9379) = \ang{3.4}$. $\vec n_1\cdot\vec
n_2 = 0.3469$, \ang{69.7} antara kedua normal, sudut antar-sayap \ang{110.3}:
sebuah lipatan terbuka.
\end{solution}

\begin{solution}{exo:b3:structural-analysis:6}
Penunjaman lebih dari \ang{60} berarti $\theta < \ang{30}$. Sama luas: $r^2 =
2\sin^2\ang{15} = 0.134$, \qty{13.4}{\percent} dari jaringnya; sama sudut:
$\tan^2\ang{15} = 0.072$, \qty{7.2}{\percent}. Pada bola, tudung itu memuat
$1 - \cos\ang{30} = 0.134$ dari belahan bumi: hanya jaring sama luas yang menampilkannya
dengan porsi sebenarnya, sehingga \omterm{def:b3:structural-analysis:stereo}{kutub} dihitung di atasnya.
\end{solution}

\begin{solution}{exo:b3:structural-analysis:7}
\begin{gather*}
  T_{xx} = \frac{2\sin^2\ang{20}}3 = 0.078,\qquad
  T_{zz} = \frac{2\cos^2\ang{20} + 1}3 = 0.922,\\
  T_{xz} = -\frac{\sin\ang{20}\cos\ang{20} - \sin\ang{20}\cos\ang{20}}3 = 0,
\end{gather*}
dan setiap suku dengan $y$ bernilai 0. Nilai eigennya adalah $0.922$ (vertikal),
$0.078$ (timur) dan $0$ (utara): \omterm{def:b3:structural-analysis:pi}{sumbu beta} horizontal, utara--selatan.
Ketiga \omterm{def:b3:structural-analysis:stereo}{kutub} terletak pada satu bidang, sehingga tidak ada yang memiliki komponen sepanjang normalnya,
dan $\vec u^{\mathsf T}T\vec u = 0$ untuk $\vec u$ sepanjang normal itu.
\end{solution}

\begin{solution}{exo:b3:structural-analysis:8}
$a = 2$, $b = 1.5$, $k = 1/0.5 = 2$: konstriksi. $Z = 1/(1.8\times1.1) =
0.505$, $a = 1.636$, $b = 2.178$, $k = 0.636/1.178 = 0.54$: penggepengan.
\end{solution}

\begin{solution}{exo:b3:structural-analysis:9}
$S_1/S_2 = 2.25$ dan $S_1S_2 = 1$: $S_1 = 1.5$, $S_2 = 0.667$; sumbu panjang
meregang sebesar \qty{50}{\percent}. Rongga itu memiliki semi-sumbu $1.5\times0.4 =
\qty{0.60}{mm}$ dan $0.667\times0.4 = \qty{0.27}{mm}$.
\end{solution}

\begin{solution}{exo:b3:structural-analysis:10}
Sama sudut: $1/(2\cos^2\ang{30}) = 0.667$ sepanjang jari-jari dan
$\tan\ang{30}/\sin\ang{60} = 0.577/0.866 = 0.667$ mengelilingi: sama, konformal.
Sama luas: $\mathrm dr/(R\,\mathrm d\theta) = (\sqrt2/2)\cos\ang{30} = 0.612$
dan $\sqrt2\sin\ang{30}/\sin\ang{60} = 0.707/0.866 = 0.816$; hasil kalinya
$0.500$, faktor luas dari teorema itu. Lingkaran kecil menjadi elips dengan
perbandingan sumbu $0.816/0.612 = 1.33$, panjang mengelilingi pusat.
\end{solution}

\begin{solution}{exo:b3:structural-analysis:11}
Sesar itu vertikal, sehingga normalnya horizontal dan tegak lurus terhadap
$\sigma_2$ yang vertikal; menurut bagian (b) teorema, gelincirnya tegak lurus terhadap
$\sigma_2$, horizontal, untuk setiap $\Phi$. Dengan $\sigma_1 = 1$ ke arah utara,
$\sigma_3 = 0$ ke arah timur: $S\vec n = (0, -0.342, 0)$, $\vec n\cdot S\vec n =
0.117$, $\vec\tau = (0, -0.342, 0) - 0.117\,(0.940, -0.342, 0) = (-0.110,
-0.302, 0)$ dan $-\vec\tau = (0.110, 0.302, 0)$, ke arah \ang{020}. Blok
di sisi $\vec n$, yaitu sisi timur, bergerak utara-utara-timur sepanjang
jurus: dilihat dari sisi barat, ia bergerak ke kiri, sebuah sesar sinistral (rake
0). $\sigma_2$ tidak masuk ke dalam perhitungan.
\end{solution}

\begin{solution}{exo:b3:structural-analysis:12}
$S = (7.5 - 6.0)/1.0 = \qty{1.5}{km}$; $h = 6.0/1.5 = \qty{4.0}{km}$ di bawah
horizon yang lebih rendah. Penampang itu panjangnya $20 + 1.5 = \qty{21.5}{km}$, dan berubah
sebesar $-1.5/21.5 = \qty{-7.0}{\percent}$.
\end{solution}

\begin{solution}{pb:b3:structural-analysis:1}
\textbf{1.} \omterm{def:b3:structural-analysis:stereo}{Kutub} berarah \ang{255} dan menunjam \ang{77}; $\theta =
\ang{13}$ dan $r = 10\sqrt2\sin\ang{6.5} = \qty{1.60}{cm}$ dari pusat,
ke arah \ang{255}.

\textbf{2.} $\vec n = (\sin\ang{13}\sin\ang{75}, \sin\ang{13}\cos\ang{75},
\cos\ang{13}) = (0.2173, 0.0582, 0.9744)$. Enam vektor satuan yang hampir sejajar
dijumlahkan menghasilkan hampir 6: $|\vec S_E| = 5.9974$.

\textbf{3.} Sayap timur: $\bar\alpha = \operatorname{atan2}(1.2583, 0.2017) =
\ang{80.9}$ dan $\bar\delta = \arccos(5.8604/5.9974) = \ang{12.3}$; $\bar R
= 5.9974/6 = 0.99956$ dan sebarannya $\sqrt{2\times0.00044} =
\qty{0.0297}{rad}$, atau \ang{1.7}. Sayap barat: $\bar\alpha =
\operatorname{atan2}(-1.2445, 0.1886) = \ang{278.6}$ dan $\bar\delta =
\arccos(5.8636/5.9972) = \ang{12.1}$; $\bar R = 0.99953$, sebaran \ang{1.8}.

\textbf{4.}
\[
  \vec S_E\times\vec S_W = (0.0774, -14.6714, 0.4883),
\]
dengan panjang $14.680$. Arah sebaliknya menunjuk ke bawah: tren adalah $\operatorname{atan2}(-0.0774, 14.6714) = \ang{359.7}$ dan
penunjaman $\arcsin(0.4883/14.680) = \ang{1.9}$.

\textbf{5.} Kosinus sudut antara normal rata-rata adalah
$32.835/35.968 = 0.9129$, yakni hasil kali titik jumlahnya dibagi
dengan panjangnya: \ang{24.1} antara kedua normal, sudut antar-sayap \ang{155.9}:
lipatan terbuka, hampir landai.

\textbf{6.}
\begin{gather*}
  \vec m_E = (0.2098, 0.0336, 0.9772),\qquad \vec m_W = (-0.2075, 0.0314, 0.9777),\\
  \vec m_E - \vec m_W = (0.4173, 0.0022, -0.0005):
\end{gather*}
normal horizontal ke arah \ang{089.7}. Bidang aksial vertikal, dengan jurus
\ang{359.7}, praktis utara--selatan: lipatannya tegak.

\textbf{7.} Sabuk itu adalah bidang yang \omterm{def:b3:structural-analysis:stereo}{kutubnya} adalah \omterm{def:b3:structural-analysis:pi}{sumbu beta}: arah kemiringan
\ang{179.7}, kemiringan $90 - 1.9 = \ang{88.1}$, hampir bidang vertikal
timur--barat. Setiap sayap memuat sumbu lipatan, sehingga setiap normal tegak lurus
terhadapnya dan terletak pada bidang yang tegak lurus padanya.

\textbf{8.} $0.04400 + 0.00144 + 0.95456 = 1.00000$. $\vec b^{\mathsf T}T\vec
b = 0.00038$, nilai eigen terkecil: $\vec b$ sampai pembulatan adalah
vektor eigen $\vec e_3$, dan sabuk yang melalui kedua \omterm{def:b3:structural-analysis:stereo}{kutub} rata-rata adalah
sabuk paling sesuai dari kedua belas data.

\textbf{9.} $\arcsin\sqrt{0.00038} = \arcsin0.0195 = \ang{1.1}$.

\textbf{10.} Sepanjang sabuk, \omterm{def:b3:structural-analysis:stereo}{kutub-kutubnya} berada sekitar \ang{12.2} dari vertikal,
di kedua sisi; $\sin^2\ang{12.2} = 0.045$, kuadrat rata-rata dari
komponen timur--barat mereka, yaitu $\lambda_2$.

\textbf{11.} Ya: \omterm{def:b3:structural-analysis:stereo}{kutub-kutub} menyimpang dari sabuk sebesar \ang{1.1} secara rata-rata, tidak
lebih dari sebaran \ang{1.7} dari pembacaan satu sayap terhadap rata-ratanya.

\textbf{12.} $\ang{1.1}/\sqrt{12} = \ang{0.3}$: penunjamannya adalah \ang{1.9} $\pm$
\ang{0.3}.

\textbf{13.} Rake sebesar \ang{-90}: sesar normal murni, gelincir kemiringan. Anderson:
$\sigma_1$ vertikal, $\sigma_3$ horizontal sepanjang arah kemiringan,
\ang{120}--\ang{300}, $\sigma_2$ horizontal sepanjang jurus,
\ang{030}--\ang{210}.

\textbf{14.} Normal dari $120/60$ terletak pada bidang vertikal yang berarah
\ang{120}, yang memuat $\vec e_1$ dan $\vec e_3$: $n_2 = 0$. Menurut bagian (b), gelincir
tegak lurus terhadap $\sigma_2$, yaitu jurus, untuk setiap $\Phi$: gelincirnya adalah
gelincir kemiringan. Sesar itu tidak memberi informasi apa pun tentang $\Phi$; dalam rumus, baik $n_2$ maupun
$\vec e_2\cdot\vec w$, sehingga pembilang dan penyebut keduanya menjadi nol.

\textbf{15.} $\vec n = (\sin\ang{55}\sin\ang{150}, \sin\ang{55}\cos\ang{150},
\cos\ang{55}) = (0.4096, -0.7094, 0.5736)$; $n_1 = 0.574$; $n_2 = 0.5\times
0.4096 + 0.866\times(-0.7094) = -0.410$.

\textbf{16.} $\vec e_1\cdot\vec w = 0.198$ dan $\vec e_2\cdot\vec w = 0.5\times
0.771 + 0.866\times0.605 = 0.909$, sehingga $\Phi = -(0.574\times0.198)/(-0.410\times
0.909) = 0.1137/0.3727 = 0.305$.

\textbf{17.} $(0.305 + 0.294 + 0.301 + 0.306)/4 = 0.30$. Rake telah
dibulatkan ke derajat terdekat; satu derajat rake menggeser $\Phi$ sekitar 0.02 pada
sesar-sesar ini, sehingga pembulatan setengah derajat menjelaskan sebaran nilainya.

\textbf{18.} $\sigma_2 - \sigma_3 = 0.3(\sigma_1 - \sigma_3)$: $\sigma_2$ berada
lebih dekat ke $\sigma_3$ daripada ke $\sigma_1$. Kedua tegangan horizontal jauh
di bawah tegangan vertikal; wilayah itu terutama meregang ke arah \ang{120}, sedikit
sepanjang \ang{030}.

\textbf{19.} Dengan $\vec n = (0.750, -0.433, 0.500)$ dan gelincir $\vec u =
(0.433, -0.250, -0.866)$, $(\vec n - \vec u)/\sqrt2$ memiliki penunjaman \ang{75} ke arah
\ang{300}: P; $(\vec n + \vec u)/\sqrt2$ memiliki penunjaman \ang{15} ke arah \ang{120}:
T. Keduanya terletak pada bidang vertikal yang sama dengan $\sigma_1$ dan $\sigma_3$, tetapi
berjarak \ang{15} dari keduanya, karena sesar membentuk sudut \ang{30}, bukan \ang{45}, dengan
$\sigma_1$.

\textbf{20.} $\tan a = \tan\ang{12.3}\times\cos(\ang{90} - \ang{80.9}) =
0.2180\times0.9874 = 0.2153$: $a = \ang{12.1}$.

\textbf{21.} $L_0 = 8000/\cos\ang{12} = \qty{8178.7}{m}$, $S =
\qty{178.7}{m}$, perubahan relatif sebesar $\cos\ang{12} - 1 =
\qty{-2.2}{\percent}$.

\textbf{22.} $4000\tan\ang{12} = \qty{850.2}{m}$; $A = \tfrac12\times8000\times
850.2 = \qty{3.40e6}{m^2}$.

\textbf{23.} $h = \num{3.40e6}/178.7 = \qty{19.0}{km}$. Detasemen di
dasar lapisan penutup akan memberi $h$ beberapa ratus meter dan luas berlebih
yang sekitar lima puluh kali lebih kecil: lipatan di sana tidak terdetasemen. Entah
basemen terlipat bersama lapisan penutup, atau sinklin tenggelam di bawah
ketinggian semula, yang membuat $A$ lebih kecil; penampang saja tidak dapat memastikannya.

\textbf{24.} Lompatan horizontal $200/\tan\ang{60} = \qty{115.5}{m}$, gelincir $200/\sin\ang{60}
= \qty{230.9}{m}$.

\textbf{25.} $\tan a = \tan\ang{60}\cos\ang{30} = 1.500$: $a = \ang{56.3}$, dan
lompatan horizontal semu adalah $200/1.5 = \qty{133.3}{m}$. Gelincir sesar
adalah ke bawah sepanjang kemiringan, ke arah \ang{120}: penampang sepanjang \ang{120} menunjukkan
panjang sebenarnya dan harus menjadi penampang yang direstorasi.

\textbf{26.} Sesar itu memotong lipatan, sehingga lebih muda dan direstorasi lebih dahulu,
menggerakkan dinding gantung kembali ke atas sesar sejauh \qty{230.9}{m}; lipatan
kemudian dibuka. \omterm{def:b3:structural-analysis:restoration}{Garis pasak} diletakkan pada bidang aksial sinklin, di
dinding bawah, tempat lipatan simetris tidak mengalami gelincir antarlapisan.

\textbf{27.} Sumbu lipatan Varden memiliki penunjaman
\textbf{\ang[text-series-to-math=true]{1.9}} ke arah utara (tren
\ang{359.7}): \ang{2} yang tertanam di wilayah itu, dipulihkan menjadi
\ang{0.3} yang dimungkinkan oleh dua belas pembacaan.
\end{solution}
